\subsection{An orientation constraint on the orthogonal Hurwitz blocks}

\begin{statement}{D23.P19}{Proved}{Actual GT acts by isometries on a fixed quadratic block}
\end{statement}
Use the degree-1344 Hurwitz orbit of D23.P17--P18. Its
degree-112 block described in D23.C23 has coordinates
\[
 w=(q_1-q_4,q_2-q_4,q_3-q_4)\in\mathbb F_3^6/\{\pm1\},
 \qquad q_i=(I-A_i)^{-1}v_i,
\]
where the four affine elements are $(v_i,A_i)$ and their
linear parts have been simultaneously transported in
$SL_2(\mathbb F_3)$ to
\[
 A_1=A_2=\begin{pmatrix}0&1\\2&1\end{pmatrix},\qquad
 A_3=A_4=\begin{pmatrix}1&1\\2&0\end{pmatrix}.
 \tag{D23.P19.1}
\]
The canonical symmetric relative cup-product form in these
coordinates is exactly
\[
 Q=\begin{pmatrix}
 0&1&1&1&1&0\\
 1&1&0&1&0&1\\
 1&0&1&1&0&2\\
 1&1&1&0&1&1\\
 1&0&0&1&1&2\\
 0&1&2&1&2&0
 \end{pmatrix}.                                            \tag{D23.P19.2}
\]
In particular, it is the same form as the independently
reconstructed D23.C23 form, with no rescaling.

Let $h\in\widehat{GT}$ have multiplier one, and suppose that
the actual permutation $V_h$ of D23.P18 preserves this block.
Then its restriction is induced by a linear map $T_h$ with
\[
                    T_h^tQT_h=Q.                          \tag{D23.P19.3}
\]
The two possible choices of $T_h$ differ by $-I$ and give the
same projective action. Thus the actual conjugator belongs
to the projective isometry group on each fixed block. The
larger projective similitude group is not an equally valid
substitute for this necessary condition.

This statement does not identify the determinant or the
spinor norm of $T_h$, and does not by itself decide
$\Delta(h)$. It does not assume compatibility with deletion
of a coordinate of an affine tuple.

\begin{proof}
Put $k=\mathbb F_3$ and $M=k^2$, with alternating form
\[
       \langle u,v\rangle=u^tJv,
       \qquad J=\begin{pmatrix}0&1\\2&0\end{pmatrix}.
\]
Let $\Pi$ be the free group on $t_1,\ldots,t_4$, and put
$t_5=(t_1t_2t_3t_4)^{-1}$. Write $\widehat\Pi$ for its full
profinite completion. The linear representation
$\rho:\widehat\Pi\to SL_2(k)$ has
$\rho(t_i)=A_i$ for $1\leq i\leq4$ and
\[
 A_5=(A_1A_2A_3A_4)^{-1}.
\]
All five $I-A_i$ are invertible. For example,
$A_1A_2A_3A_4=\left(\begin{smallmatrix}2&1\\0&2\end{smallmatrix}\right)$,
so this holds at the fifth puncture as well.

\paragraph{Affine lifts are cohomology classes.}
An affine lift of $\rho$ is a map
$g\mapsto(c(g),\rho(g))$ into $M\rtimes SL_2(k)$, where
\[
                 c(gh)=c(g)+\rho(g)c(h).                  \tag{D23.P19.4}
\]
Thus it is a left 1-cocycle. Its four generator values are
arbitrary; the universal property of the free profinite group
extends any such assignment continuously, since the target
semidirect product is finite. Conjugating the affine lift by
a translation changes all its fixed points $q_i$ by the same
vector. Consequently the quotient by translations is
$H^1(\widehat\Pi,M)$, of dimension $8-2=6$, and setting
$q_4=0$ gives precisely the difference coordinates in the
statement. The simultaneous centralizer of (D23.P19.1) in
$SL_2(k)$ is $\{\pm I\}$. Quotienting further by this
centralizer gives the projective nonzero classes and the zero
class. D23.C23 identifies the 112 points of the indicated
Hurwitz block with the singular nonzero projective classes.

For clarity, this description concerns the fiber of the
linear-part representation class. In D23.C23 every point in
the indicated block has that linear-part class. Its other
translates are the other eleven blocks; the ordinary Hurwitz
action commutes with projection to linear parts. Thus fixing
one of these blocks fixes the corresponding linear
representation class up to $SL_2(k)$ conjugacy.

\paragraph{Relative cohomology and its orientation.}
Use the five peripheral subgroups
$C_i=\overline{\langle t_i\rangle}\cong\widehat{\mathbb Z}$
and the relative continuous cochain complex, the shifted
mapping cone of
\[
 C^\bullet(\widehat\Pi,M)\longrightarrow
              \bigoplus_{i=1}^5 C^\bullet(C_i,M).
\]
Because $I-A_i$ is invertible, both
$H^0(C_i,M)$ and $H^1(C_i,M)$ vanish. Indeed these groups are
respectively $\ker(A_i-I)$ and $M/(A_i-I)M$. The long exact
sequence therefore identifies relative $H^1$ with ordinary
$H^1(\widehat\Pi,M)$. This is the usual parabolic situation,
with no local invariant vectors.

For trivial coefficients $k$, the same sequence gives
\[
 H^2(\widehat\Pi,\{C_i\};k)
   =\operatorname{coker}\left(
       k^4\longrightarrow k^5\right)\cong k,
 \quad
 (a_1,\ldots,a_4)\longmapsto
       (a_1,\ldots,a_4,-\textstyle\sum_i a_i).              \tag{D23.P19.5}
\]
Here $H^2(\widehat\Pi,k)=0$: a continuous finite-valued
2-cocycle factors through a finite quotient, and the
corresponding finite extension splits after pullback to the
free profinite group, by its projectivity. The quotient in
(D23.P19.5) is oriented by the sum of the five coordinates.

Suppose an automorphism $\alpha$ of $\widehat\Pi$ carries
each $t_i$ to a conjugate of $t_i^\lambda$, for one
$\lambda\in\widehat{\mathbb Z}^{\times}$. It acts by
$\lambda\bmod3$ on each of the five trivial-coefficient
peripheral $H^1$ groups, since conjugation acts trivially on
them. Consequently its action on (D23.P19.5) is also
$\lambda\bmod3$. This proves the orientation character
directly, without realizing $\alpha$ by a homeomorphism.
The cochain maps for conjugated peripheral subgroups are
defined using their usual conjugation identifications; these
are compatible with restriction and cup product. Changing
one of the conjugating elements gives the corresponding
inner cochain homotopy and does not change the cohomology
map or this orientation character.

Relative cup product, followed by the alternating coefficient
pairing $J$, therefore gives a symmetric bilinear pairing
\[
 H^1(\widehat\Pi,M)\times H^1(\widehat\Pi,M)
       \longrightarrow H^2(\widehat\Pi,\{C_i\};k)\cong k.
 \tag{D23.P19.6}
\]
The symmetry is the product of the degree-one cup sign and
the alternating coefficient sign. Its naturality shows that
pullback by $\alpha^{-1}$ multiplies it by
$\lambda^{-1}\bmod3$, when the coefficient forms are
identified isometrically.

There is no unproved discrete-to-profinite comparison in
this calculation. Restriction to the dense discrete free
group identifies $H^0$ and $H^1$ for these finite modules:
the same arbitrary generator-value description proves the
claim in degree one. It identifies the corresponding
peripheral groups in degrees zero and one. The discrete
free group also has $H^2=0$, so (D23.P19.5) is carried
isomorphically to its discrete counterpart. Restriction of
cochains respects cup product. These observations suffice
for all the comparison maps used in (D23.P19.6); no
comparison in higher degrees is required.

\paragraph{Exact cup-product matrix.}
Here is an explicit bar-chain computation, fixing the sign
and the coordinate identification. For two cocycles $c,d$
put
\[
 c_i=c(t_i),\quad d_i=d(t_i),\quad
 c_i=(I-A_i)q_i,\quad
 p_i=t_1\cdots t_i,\quad P_i=A_1\cdots A_i.
\]
The fifth value is determined by $c(p_5)=0$. In normalized
bar chains the chain
\[
                  Z=\sum_{i=1}^5[p_{i-1}\mid t_i]
\]
has boundary $\sum_i[t_i]$. Thus
$(Z,-\sum_i[t_i])$ is a relative cycle. A relative lift of
$c$ has boundary primitives $-q_i$, because
$\delta(-q_i)(t_i)=(I-A_i)q_i=c_i$. The ordinary cup
cochain is
\[
               (g,h)\longmapsto
                 \langle c(g),\rho(g)d(h)\rangle.
\]
Its relative boundary component is
$\langle-q_i,d_i\rangle$. Evaluation on the displayed
relative cycle consequently gives
\[
 B(c,d)=\sum_{i=1}^5
       \left(\langle c(p_{i-1}),P_{i-1}d_i\rangle
                       +\langle q_i,d_i\rangle\right).
 \tag{D23.P19.7}
\]
This also independently fixes the orientation convention in
the matrix calculation. Evaluating (D23.P19.7) on the eight
fixed-point coordinate vectors $(q_1,q_2,q_3,q_4)$ gives a
symmetric matrix of rank six whose radical is exactly the
two-dimensional diagonal translation subspace. On the
section $q_4=0$ its matrix is exactly (D23.P19.2).
The certificate evaluates all 36 entries, checks the full
eight-dimensional radical, and compares the resulting
matrix literally with D23.C23. In particular the form is
nondegenerate on $H^1$, and no appeal to invariant-form
uniqueness is needed for its identification.

\paragraph{Application to actual GT.}
By D23.P18, the actual permutation is induced by
precomposition with $\kappa_h^{-1}$, where $\kappa_h$ is the
restriction of the Drinfeld automorphism to the free kernel
of the rightmost forgetting map. That proof gives the
required peripheral powers for $t_1,\ldots,t_4$. It also
gives the fifth: the boundary product is
$t_1t_2t_3t_4=z_4^{-1}z_5$, and the standard action sends
each full twist $z_j$ to $z_j^\lambda$. The two nested full
twists commute, so the boundary product, and hence its
inverse $t_5$, has the same power $\lambda$.

If the block is fixed, choose $B\in SL_2(k)$ with
\[
 B\rho(\kappa_h^{-1}(g))B^{-1}=\rho(g)
                  \qquad(g\in\widehat\Pi).
\]
The induced map on affine lifts is
\[
                     c(g)\longmapsto
                       Bc(\kappa_h^{-1}(g)).              \tag{D23.P19.8}
\]
It is linear on cocycles, preserves coboundaries, and is
bijective. Since $B\in SL_2(k)$, it preserves $J$.
The naturality calculation therefore gives
\[
                B(T_hc,T_hd)=\lambda^{-1}B(c,d).
\]
For multiplier one this is (D23.P19.3). The only ambiguity
in $B$ is the simultaneous centralizer $\{\pm I\}$, which
changes $T_h$ by $-I$ and has no effect on its projective
action or on its multiplier. This proves the statement.
\end{proof}

\paragraph{Source and check scope.}
The parabolic cup pairing and the reversal of symmetry type
are standard; a primary reference is Dettweiler--Wewers,
\emph{Variation of parabolic cohomology and Poincar\'e
duality}, arXiv:math/0411119v1, Corollary 2.2 and \S2.3,
printed p.10. Their Theorem 2.5, printed p.12, supplies an
explicit formula in their right-cocycle convention:
\url{https://arxiv.org/pdf/math/0411119}.
Here the needed orientation character, the comparison in the
free profinite case, and formula (D23.P19.7) in the actual
left-cocycle marking were proved directly above.
The exact full-twist power normalization is also stated
explicitly in Lochak--Nakamura--Schneps,
arXiv:2602.12462v1, \S2, Theorem 2.1, printed pp.2--3:
$F(\omega_i)=\omega_i^\lambda$ with
$\omega_i=(\sigma_1\cdots\sigma_{i-1})^i$.
\url{https://arxiv.org/pdf/2602.12462}.

The exact certificate is
\path{hurwitz_cup_pairing.py}; its record is
\path{hurwitz_cup_pairing.json}. It identifies the matrix and
checks all thirteen supplied block-stabilizer matrices
against it. It does not purport to enumerate, approximate,
or construct an unknown full-GT element.
